\section{BAB V: KESIMPULAN DAN SARAN}

\subsection{Kesimpulan}
Berdasarkan perancangan, implementasi, dan serangkaian pengujian yang telah dilakukan pada Proyek 3 sistem GaluraSwara, dapat diambil beberapa kesimpulan:
\begin{enumerate}
    \item Ekosistem aplikasi GaluraSwara berhasil dikembangkan secara menyeluruh yang mengintegrasikan aplikasi seluler Native Kotlin, layanan backend Gocroot Go Fiber, dokumentasi OpenAPI 3.0, serta landing page Croot.js.
    \item Pemutar musik berbasis AndroidX Media3 (ExoPlayer) pada aplikasi Android terbukti andal dalam menyajikan pemutaran lagu tradisional berkualitas tinggi dengan transisi status yang responsif pada antarmuka Jetpack Compose.
    \item Penerapan otomasi DevOps menggunakan GitHub Actions berhasil menyederhanakan siklus rilis; pembaruan kode backend secara otomatis dideploy ke Alwaysdata, dan rilis APK Android secara mandiri di-build, diberi tag versi (\textit{autotag}), serta dipublikasikan ke GitHub Releases.
    \item Desain bernuansa etnik tradisional pada seluruh media antarmuka memberikan ciri khas identitas budaya Nusantara yang kuat dan ramah pengguna.
\end{enumerate}

\subsection{Saran}
Untuk pengembangan dan penyempurnaan sistem GaluraSwara di masa mendatang, disarankan beberapa hal sebagai berikut:
\begin{enumerate}
    \item Menambahkan fitur pengunduhan lagu secara lokal (\textit{offline caching}) pada aplikasi Android agar pengguna dapat mendengarkan musik tradisional tanpa koneksi internet.
    \item Menambahkan fitur lirik lagu berbasis waktu (\textit{synchronized lyrics / rumpaka}) untuk memudahkan pengguna mempelajari lirik lagu-lagu daerah Sunda dan Nusantara.
    \item Mengintegrasikan analitik pendengar pada backend untuk memetakan jenis kesenian tradisional yang paling digemari oleh masyarakat.
\end{enumerate}
